\subsection*{MISC Library}

The following
names are provided by the MISC library:
\begin{itemize}
  \item
    \verb#time_time#\textt{()}: \textit{primitive},
    returns number of milliseconds elapsed since
    January 1, 1970 00:00:00 UTC
  \item \verb#int#\texttt{(x)}, \verb#float#\texttt{(x)},
    \verb#bool#\texttt{(x)}, \verb#complex#\texttt{(x)}, \verb#str#\texttt{(x)}: \textit{primitive},
    converts the value \texttt{x} to an integer, a float, a boolean, a complex number, or a string
    representation of \texttt{x} respectively, and returns the
    respective value. For \verb#int#\texttt{(x)}, if \texttt{x} is a string,
    \texttt{x} is parsed
    as an integer in base 10.
  \item \verb#int#\texttt{(s, i)}: \textit{primitive},
    interprets the \emph{string} \texttt{s} as an integer, using the
    positive integer \texttt{i} as radix, and returns the respective value.
  \item \verb#repr#\texttt{(s)}: \textit{primitive},
    returns a string representation of \texttt{s} that is unambiguous and suitable for
    debugging; for many values, this string can be used to reconstruct \texttt{s}. 
  \item
    \href{https://sourceacademy.org/sicpjs/4.1.2\#p2}{\lstinline{is_boolean(x)}},
    \href{https://sourceacademy.org/sicpjs/2.3.2\#p5}{\lstinline{is_int(x)}},
    \href{https://sourceacademy.org/sicpjs/2.3.2\#p5}{\lstinline{is_float(x)}},
    \href{https://sourceacademy.org/sicpjs/2.3.2\#p7}{\lstinline{is_string(x)}},
    \href{https://sourceacademy.org/sicpjs/4.1.2\#p2}{\lstinline{is_none(x)}},
    \verb#is_function#\texttt{(x)},
    \verb#is_complex#\texttt{(x)}: \textit{primitive}, returns
    \texttt{True} if the type of \texttt{x} matches the function name
    and \texttt{False} if it does not.
  \item \texttt{real(x)}, \texttt{imag(x)}: \textit{primitive}, return the 
    real and imaginary parts of a complex number \texttt{x} respectively.
  \item \texttt{input(s)}: \textit{primitive}, pops up a window that
    displays the \emph{string} \texttt{s}, provides
    an input line for the user to enter a text, a ``Cancel'' button
    and an ``OK'' button. The call of \texttt{input}
    suspends execution of the program until one of the two buttons is
    pressed. If
    the ``OK'' button is pressed, \texttt{input} returns the entered
    text as a string.
    If the ``Cancel'' button is pressed, \texttt{input} returns a
    non-string value.
  \item
    \verb#print#\texttt{(*x)}:
    \textit{primitive}, takes in any number of parameters, converts them to their string
    representations using \verb#str# and writes them to the console, followed by a newline character.
  \item
    \verb#error#\texttt{(*x)}:
    \textit{primitive}, takes in any number of parameters, converts them to their string
    representations using \verb#str# and writes them to the console, followed by a newline character.
    The evaluation of any call of \texttt{error} aborts the running program immediately.
\end{itemize}
All library functions can be assumed to run
in $O(1)$ time, except \texttt{print}, \texttt{error}, \texttt{str} and \texttt{repr},
which run in $O(n)$ time, where $n$ is
the size (number of components such as pairs) of their first argument.
